\section{2026 年 11 月核验增补}

本节按论文首次提交时间排列。每项方法附有可核验的外部入选依据。

% METHOD: 2603.09703
\paragraph{ProGS：把高斯场编码为逐层可解码的单一码流}
2026.03~\cite{2603_09703}。普通高斯压缩生成完整场景的码流，接收不完整时未必能渲染。ProGS 从 SfM 点云建立八叉树锚点，要求每个细层锚点依赖已解码的父层。编码后的共享头随后是多个累计前缀，接收一层即可解码并显示；之后追加层增加细节。低层的轻量残差用于改善早期可用画面，同时保留基础拓扑。

原文在 Mip-NeRF 360、BungeeNeRF 和 Tanks and Temples 比较码率与质量，区分早期前缀和完整码流。前缀体积需要计入共享头和已收块，不能用训练 PLY 大小替代；这组结果也不能直接推出实际网络上的端到端缓冲时延。官方仓库有训练脚本和独立编码、解码、渲染入口。

\textbf{入选依据} \href{https://github.com/zhiyetang/progs-official}{公开实现}已找到与论文对应的代码文件。

% METHOD: 2607.05667
\paragraph{CGVQ：把二维图像高斯的属性分组后量化}
2026.07~\cite{2607_05667}。CGVQ 将单张图像建模为二维高斯集合，原语包括二维位置、旋转尺度矩阵及颜色系数。方法利用属性聚类分配码本，并以量化后的原语渲染图像，利用同类属性的统计相似性减少码率。

作者在 Kodak 图像上用每像素比特与 PSNR、SSIM 比较 GaussianImage，并说明实验选择的固定聚类数量适合文中的原语规模。该任务是二维图像压缩，不可用来断言多视图三维高斯场景的率失真收益。作者单位为 NVIDIA。

\textbf{入选依据} 论文署名单位包括 NVIDIA（\href{https://arxiv.org/html/2607.05667v1}{论文原文}）。

% METHOD: 2607.12656
\paragraph{SpeedyGS：先形成内容自适应结构，再选择编码复杂度}
2026.07~\cite{2607_12656}。SpeedyGS 先以率失真代理训练高斯表示，通过按内容调整的量化和剪枝减少原语，再用可选复杂度的统计编码压缩几何与属性。学习式编码用稀疏八叉树几何符号及一维属性符号；更轻配置服务不同成本条件。

在论文的 Mip-NeRF360 对照中，完整方案把基线的 744.70 MB 降到 6.51 MB，渲染从 124 提至 208 FPS；轻量配置则采用不同压缩率。结果是选定数据、模型和计时方式下的有损率失真对照。作者单位包括 Huawei Technologies。

\textbf{入选依据} 论文署名单位包括 Huawei（\href{https://arxiv.org/html/2607.12656v1}{论文原文}）。

% METHOD: 2607.14513
\paragraph{GPU 纹理压缩：将已有高斯的球谐颜色属性写入块编码}
2026.07~\cite{2607_14513}。该方法将高斯的球谐颜色系数组合为适用现有 GPU 图形纹理格式的块，利用固定大小、硬件可读的块格式减少属性存储。不同系数打包方式与纹理格式影响码字大小和渲染失真。

论文只压缩三组场景 bicycle、bonsai 和 garden 的球谐部分，原语位置、形状和不透明度仍使用 32 位浮点；用 BC1 与 BC7 等配置比较球谐码流和 PSNR。表中每原语球谐字节不能理解为完整的场景存储或吞吐收益。作者单位包括 Qualcomm AI Research 和 Qualcomm Graphics Research。

\textbf{入选依据} 论文署名单位包括 Qualcomm（\href{https://arxiv.org/html/2607.14513v1}{论文原文}）。

% METHOD: 2607.22956
\paragraph{ParticleGS：以图像保真目标压缩科学粒子可视化}
2026.07~\cite{2607_22956}。ParticleGS 直接从科学粒子数据训练紧凑高斯可视化代理，以渲染图像而非单粒子误差作为优化目标。VizMapper 按显示参数调整外观，空间分块让同一思路适用于大规模粒子场，查看时无需逐次解压原始粒子。

作者在 HACC 与 FIRE-2 场景比较视觉失真、存储和渲染速度。高斯混合模型只能近似回采粒子并保留部分大尺度密度统计，不能逐点还原原数据。官方仓库列有从原始粒子训练、渲染和比对指标的复现流程。

\textbf{入选依据} \href{https://github.com/bojiang03/particlegs}{公开实现}已找到与论文对应的代码文件。

% METHOD: 2607.28049
\paragraph{TSOG：为已有四维高斯内容定义可交换的存储格式}
2026.07~\cite{2607_28049}。TSOG 在面向静态高斯属性图像的 SOG 上加入共享时间轴、原语生存周期和几何外观的时间参数。各属性图像中同一个高斯仍对应同一个像素位置，因此可以在旧表示基础上加入动态内容。它规定如何存储已有四维高斯，不重新学习四维重建模型。

实验用体积视频挑战数据及 Neural 3D 的既有输出比较文件大小和图像质量，作者报告相对 PLY 体积减少超过九成。该结果依赖实验编码口径，持续流送和更复杂时间模型仍待评估。原文署名单位列有 Xiaomi Technology Netherlands。

\textbf{入选依据} 论文署名单位包括 Xiaomi（\href{https://arxiv.org/html/2607.28049v1}{论文原文}）。

% METHOD: 2609.02184
\paragraph{CC-4DGS：分别压缩动态变形和规范高斯属性}
2026.09~\cite{2609_02184}。CC-4DGS 用计算式变形场替换多分辨率哈希表，并压缩规范坐标中的高斯属性。该组合同时针对动态网络存储与点云存储，便于在保留动态渲染的前提下减少模型容量。

作者在 N3DV 和 Technicolor 与多种动态方法比较，并对 Swift4D 和本方法施加相同点云压缩以对齐存储比较。Swift4D 仍具有较高渲染速率，因此存储收益和实时速度属于需要分别衡量的量。

\textbf{入选依据} \href{https://github.com/kyungdaepark/cc-4dgs}{公开实现}已找到与论文对应的代码文件。

% METHOD: 2609.03334
\paragraph{分频率阶段训练多个高斯场以降低活跃原语数量}
2026.09~\cite{2609_03334}。该方法先将图像拆成低频底层和拉普拉斯残差，再从粗到细训练各频带的高斯场。训练时把不再活跃的层归档，以减少同时参与优化的原语；推断时分别渲染各层并合成结果。

论文在三类数据集的 1K 设置比较训练时间、原语量和画质。训练活跃内存的减少不能直接解释为推断只需一个场，部署仍需统计各层合成的成本。

\textbf{入选依据} 论文署名单位包括 XGRIDS（\href{https://arxiv.org/abs/2609.03334}{论文原文}）。\href{https://github.com/sorenzhang574/laplacian-gs}{公开实现}已找到与论文对应的代码文件。

% METHOD: 2609.30245
\paragraph{实用高斯压缩：在不依赖空间聚合的条件下整数化解码}
2026.09~\cite{2609_30245}。这项工作由坐标派生几何上下文，再对紧凑锚点潜变量与属性作因果编码。线性变换与激活构成的上下文网络配合率失真学习、量化感知训练和整数推理，简化解码并降低平台浮点差异。

论文在多个公开场景比较码流大小、渲染质量和总编码、解码时间，区分不同训练预算与量化感知配置。跨平台一致性是实验所用整数解码路径的结果，不能据此对未测硬件或任意码流作兼容性保证。

\textbf{入选依据} \href{https://github.com/pengpeng-yu/cosa-gs}{公开实现}已找到与论文对应的代码文件。

% METHOD: 2610.05576
\paragraph{SteadySplats：以空间复用和方差训练降低随机渲染噪声}
2026.10~\cite{2610_05576}。SteadySplats 将低样本随机高斯渲染的噪声分为估计方差与样本数问题。一方面复用邻域和历时样本，另一方面训练时约束同一视线上贡献颜色的方差，使少量抽样也能接近按序透明度合成的结果。

论文以 7 个 Mip-NeRF 360 场景的 21 条相机轨迹、3152 帧验证，并将按序混合渲染作为参照。作者指出方差训练会稀释光度损失、略降充分收敛时的部分峰值质量指标；样本复用也并非零开销。

\textbf{入选依据} 论文署名单位包括 Google DeepMind、NVIDIA（\href{https://arxiv.org/abs/2610.05576}{论文原文}）。

% METHOD: 2610.07795
\paragraph{GSCV：用跨帧相容的属性图接入标准视频编码器}
2026.10~\cite{2610_07795}。GSCV 以参考布局约束相邻帧的属性排列，使同一三维结构在二维高位深属性图中尽量能被帧间预测跟踪。属性图再由视频编码器压缩，以利用已有标准而保持较高的精度上限。

论文比较 MPEG 测试中的已追踪和半追踪高斯序列，以 GSCodec Studio 与 G-PCC 为对照，报告码率和重建质量。作者指出重排后的属性图不同于自然视频，通用编码器仍会限制压缩率；多个视频属性流也会改变系统整合开销。

\textbf{入选依据} 论文署名单位包括 Qualcomm（\href{https://arxiv.org/html/2610.07795v1}{论文原文}）。
